\section{Appendix: Identification and upper-bound proofs}
\label{sec:upper-appendix}

\subsection{Identification from primitive variables}

The causal interpretation of the transported contrasts rests on consistency, conditional instrument randomization, exclusion, and monotonicity. The two transport restrictions then carry the source conditional contrasts to target complier quantities.

\begin{lemmav}{lem:transported-cace-identification}[Transported complier effect identification]
Let \(n\in\mathbb N\) and let \(P\) be a transport law satisfying:
\begin{itemize}
\item \textbf{(Lower envelope.)} \(0<c_f<1\), where \(c_f\) is the lower density envelope.
\item \textbf{(Upper envelope.)} \(1<C_f\), where \(C_f\) is the upper density envelope.
\item \textbf{(Smoothness radius.)} \(1<L\), where \(L\) is the common Hölder radius.
\item \textbf{(Model membership.)} \(P\in\mathcal M_n\), the model class of \cref{obj:def:model-class} with constants \(c_f,C_f,L\).
\end{itemize}
Write \(P_n^F\) for the full-data law, \(D(z)\) for potential receipt under encouragement \(z\), and \(Y(d)\) for potential outcome under receipt \(d\). For each \(z\in\{0,1\}\), there are versions of the source conditional means such that, almost everywhere under the source covariate law,
\[
m_{Dz}(x)
 =\mathbb E_{P_n^F}[D(z)\mid X=x,S=1],
\qquad
m_{Yz}(x)
 =\mathbb E_{P_n^F}[Y(D(z))\mid X=x,S=1].
\]
Here \(m_{Dz}\) and \(m_{Yz}\) are the source receipt and outcome arm means.

Define the source arm contrasts and the complier indicator by
\[
\Delta_D(x)=m_{D1}(x)-m_{D0}(x),
\qquad
\Delta_Y(x)=m_{Y1}(x)-m_{Y0}(x),
\qquad
C=\mathbf 1\{D(1)=1,\ D(0)=0\}.
\]
There are common versions of the conditional means in each of the following chains for which, for \(P_{\mathrm T}\)-almost every \(x\), where \(P_{\mathrm T}\) is the target covariate law,
\[
\begin{aligned}
\Delta_D(x)
&=\mathbb E_{P_n^F}[D(1)-D(0)\mid X=x,S=1]\\
&=\mathbb E_{P_n^F}[C\mid X=x,S=1]
 =\mathbb E_{P_n^F}[C\mid X=x,S=0],\\
\Delta_Y(x)
&=\mathbb E_{P_n^F}[Y(D(1))-Y(D(0))\mid X=x,S=1]\\
&=\mathbb E_{P_n^F}[(Y(1)-Y(0))C\mid X=x,S=1]\\
&=\mathbb E_{P_n^F}[(Y(1)-Y(0))C\mid X=x,S=0].
\end{aligned}
\]
Moreover,
\[
D(1)-D(0)=C
\qquad P_n^F\text{-almost surely}.
\]
For \(A\in\{D,Y\}\), define the transported contrast using the target covariate density \(f_{\mathrm T}\) by
\[
T_A=\int \Delta_A(x)f_{\mathrm T}(x)\,dx.
\]
The transported first stage equals the target complier share, and the transported outcome contrast equals the target complier-weighted outcome effect:
\[
T_D=\mu=P_n^F(C=1\mid S=0),
\qquad
T_Y=\mathbb E_{P_n^F}[(Y(1)-Y(0))C\mid S=0].
\]
The target complier average treatment effect \(\theta\) satisfies
\[
\theta
=\frac{\mathbb E_{P_n^F}[(Y(1)-Y(0))C\mid S=0]}
       {P_n^F(C=1\mid S=0)}
=\frac{T_Y}{T_D}
\in\Theta,
\]
where \(\Theta\) is the effect domain specified by the model.
\end{lemmav}

The common-version qualification makes the conditional identities compatible with target integration. Positive density envelopes provide common covariate support, monotonicity identifies the receipt contrast with complier prevalence, and exclusion identifies the outcome contrast with the complier-weighted receipt effect. Positive transported strength makes their ratio the target complier average effect.

\subsection{Coordinate and sampling conventions}

The estimator uses one target-density coordinate and six marked source coordinates. The next two definitions fix the coordinate names, observable marks, channels, and index conversion once for the remaining construction.

\begin{definitionv}{synth_3}[Exact block size]
The block size \(m_b\), for a nonnegative integer \(n\) and an index \(b\in\{0,1,2,3\}\), is defined by
\[
m_b := \max\left\{
\left\lfloor \frac{(b+1)n}{4} \right\rfloor
-
\left\lfloor \frac{bn}{4} \right\rfloor,
0
\right\}.
\]
\end{definitionv}

\begin{definitionv}{synth_2}[Deterministic sample blocks]
The deterministic block \(\mathcal B_b^G\) in sample channel \(G\), for sample size \(n\in\{0,1,2,\ldots\}\) and block index \(b\in\{0,1,2,3\}\), uses zero-based record indices and is defined by
\[
\mathcal B_b^G
=
\left\{
i\in\mathbb Z:
0\le i<n,\quad
\left\lfloor\frac{bn}{4}\right\rfloor
\le i<
\left\lfloor\frac{(b+1)n}{4}\right\rfloor
\right\}.
\]
\end{definitionv}

\begin{definitionv}{synth_19}[Target-density coordinate]
In the marked-density vector, \(t\) denotes the target covariate density:
\[
t:=f_{\mathrm T},\qquad t(x)=f_{\mathrm T}(x),\quad x\in[0,1].
\]
Thus \(F_0=t=f_{\mathrm T}\), and the named-coordinate subscript \(t\) denotes coordinate \(0\): \(V_t:=V_0\) for \(V=F\), \(V=H_K^{(b)}\), \(V=R_K^{(b)}\), or \(V=W_r\). In particular, \(F_t=f_{\mathrm T}\). In \(\operatorname{bump}(t)\) and \(\int_{\Theta}P(t\in C_n)\,dt\), \(t\) is instead a local real argument or integration variable, respectively.
\end{definitionv}

\begin{definitionv}{synth_14}[Coordinate marks \(W_{i,r}\) and sample blocks \(\mathcal B_b^G\)]
The coordinate index set and coordinate ordering are
\[
\mathcal I=\{0,\ldots,6\},
\qquad
(F_0,F_1,F_2,F_3,F_4,F_5,F_6)
=(t,q_0,q_1,r_{Y0},r_{Y1},r_{D0},r_{D1}).
\]
The sample channel of coordinate \(i\in\mathcal I\) is
\[
g(i)=
\begin{cases}
\mathrm T,&i=0,\\
\mathrm S,&i\in\{1,\ldots,6\}.
\end{cases}
\]
Vector, histogram, residual, and mark subscripts follow the correspondence
\[
\begin{array}{c|ccccccc}
i&0&1&2&3&4&5&6\\ \hline
\text{coordinate}&t&q_0&q_1&r_{Y0}&r_{Y1}&r_{D0}&r_{D1}.
\end{array}
\]
For \(V=F\), \(V=H_K^{(b)}\), \(V=R_K^{(b)}\), or \(V=W_r\), the named-coordinate subscripts satisfy
\[
(V_t,V_{q_0},V_{q_1},V_{r_{Y0}},V_{r_{Y1}},V_{r_{D0}},V_{r_{D1}})
=(V_0,V_1,V_2,V_3,V_4,V_5,V_6).
\]

Index the source records and target covariate records from zero:
\[
O_r^{\mathrm S}=(X_r^{\mathrm S},Z_r,D_r,Y_r),
\qquad
X_r^{\mathrm T},
\qquad r=0,\ldots,n-1.
\]
Let \(Y_r^{\mathrm{clip}}\) denote the clipped observable outcome used in
\cref{obj:synth_7,obj:synth_8,obj:def:cubic-estimator}.
The observable coordinate marks, with this same clipping convention throughout, are
\[
\begin{aligned}
W_{0,r}&=1,\\
W_{1,r}&=\mathbf 1\{Z_r=0\},
&
W_{2,r}&=\mathbf 1\{Z_r=1\},\\
W_{3,r}&=\mathbf 1\{Z_r=0\}Y_r^{\mathrm{clip}},
&
W_{4,r}&=\mathbf 1\{Z_r=1\}Y_r^{\mathrm{clip}},\\
W_{5,r}&=\mathbf 1\{Z_r=0\}D_r,
&
W_{6,r}&=\mathbf 1\{Z_r=1\}D_r.
\end{aligned}
\]
The mark \(W_{0,r}\) accompanies target record \(r\); the remaining marks accompany source record \(r\).

For \(G\in\{\mathrm S,\mathrm T\}\) and \(b\in\{0,1,2,3\}\), the blocks and their sizes are
\[
\mathcal B_b^G
=
\left\{
r\in\{0,\ldots,n-1\}:
\left\lfloor\frac{bn}{4}\right\rfloor
\le r<
\left\lfloor\frac{(b+1)n}{4}\right\rfloor
\right\},
\qquad
m_b=\left\lfloor\frac{(b+1)n}{4}\right\rfloor
-\left\lfloor\frac{bn}{4}\right\rfloor.
\]
Their one-based counterparts are
\[
\widetilde{\mathcal B}_b^G
=\{r+1:r\in\mathcal B_b^G\}
=\{\lfloor bn/4\rfloor+1,\ldots,\lfloor(b+1)n/4\rfloor\}.
\]
One-based covariates and marks are defined together by
\[
\widetilde X_s^G=X_{s-1}^G,
\qquad
\widetilde W_{i,s}=W_{i,s-1},
\qquad s=1,\ldots,n.
\]
Thus, for a covariate set \(E\), the corresponding block sums satisfy
\[
\sum_{s\in\widetilde{\mathcal B}_b^{g(i)}}
\widetilde W_{i,s}\,
\mathbf 1\{\widetilde X_s^{g(i)}\in E\}
=
\sum_{r\in\mathcal B_b^{g(i)}}
W_{i,r}\,\mathbf 1\{X_r^{g(i)}\in E\}.
\]
Block \(0\) supplies the training records, and blocks \(1,2,3\) supply the evaluation records. Each coordinate-\(i\) sum uses covariates and marks with the same record index in channel \(g(i)\).
\end{definitionv}

\subsection{Histograms and cubic estimator}

A normalized marked histogram estimates each coordinate on a common dyadic partition. Block 0 trains the clipped pilot; blocks 1--3 provide independent factors for the correction terms.

\begin{definitionv}{synth_7}[Blockwise marked histogram definition]
The marked histogram \(H_{i,K}^{(b)}(x)\), for coordinate \(i\), resolution \(K\), block \(b\), and evaluation point \(x\), is defined by
\[
H_{i,K}^{(b)}(x)=
\begin{cases}
0,&n<256,\\[2mm]
\displaystyle
\frac{K}{m_b}
\sum_{\ell=0}^{K-1}
\mathbf 1\{x\in E_{\ell,K}\}
\sum_{r\in\mathcal B_b^{g(i)}}
W_{i,r}\,
\mathbf 1\{X_r^{g(i)}\in E_{\ell,K}\},
&n\ge256.
\end{cases}
\]
Here \(\mathcal S_n\), the sample of source records, and \(\mathcal T_n\), the sample of target covariates, each contain \(n\) records, and the parameters range over
\[
n,K\in\{0,1,\ldots\},\qquad
i\in\mathcal I=\{0,\ldots,6\},\qquad
b\in\{0,\ldots,3\},\qquad
x\in\mathbb R.
\]
For \(K\ge1\), the cells \(E_{\ell,K}\) are
\[
E_{\ell,K}=
\begin{cases}
[\ell/K,(\ell+1)/K),&0\le\ell<K-1,\\
[\ell/K,1],&\ell=K-1.
\end{cases}
\]
For \(K=0\), the cell sum is empty and \(H_{i,0}^{(b)}(x)=0\).

Both samples use zero-based record indices \(r\in\{0,\ldots,n-1\}\), with an empty index set when \(n=0\). For sample channel \(G\in\{\mathrm S,\mathrm T\}\), the deterministic block \(\mathcal B_b^G\) and its size \(m_b\) are
\[
\mathcal B_b^G
=
\left\{
r\in\{0,\ldots,n-1\}:
\left\lfloor\frac{bn}{4}\right\rfloor
\le r<
\left\lfloor\frac{(b+1)n}{4}\right\rfloor
\right\},
\qquad
m_b=
\left\lfloor\frac{(b+1)n}{4}\right\rfloor
-
\left\lfloor\frac{bn}{4}\right\rfloor.
\]
The sample channel \(g(i)\) assigned to coordinate \(i\) is
\[
g(i)=
\begin{cases}
\mathrm T,&i=0,\\
\mathrm S,&i\in\{1,\ldots,6\}.
\end{cases}
\]
The covariate \(X_r^{g(i)}\) is the observed covariate of zero-based record \(r\) in that channel.

For zero-based source record \(r\), let \(Z_r\in\{0,1\}\) denote encouragement, \(D_r\in\{0,1\}\) treatment receipt, and \(Y_r\in\mathbb R\) the observed outcome. The observable marks \(W_{i,r}\) are
\[
\begin{aligned}
W_{0,r}&=1,\\
W_{1,r}&=\mathbf 1\{Z_r=0\},
&
W_{2,r}&=\mathbf 1\{Z_r=1\},\\
W_{3,r}&=\mathbf 1\{Z_r=0\}\max\{0,\min\{1,Y_r\}\},
&
W_{4,r}&=\mathbf 1\{Z_r=1\}\max\{0,\min\{1,Y_r\}\},\\
W_{5,r}&=\mathbf 1\{Z_r=0\}\mathbf 1\{D_r=1\},
&
W_{6,r}&=\mathbf 1\{Z_r=1\}\mathbf 1\{D_r=1\}.
\end{aligned}
\]
The constant mark \(W_{0,r}\) accompanies zero-based target record \(r\); each remaining mark accompanies zero-based source record \(r\). In every block sum, the mark and covariate use the same record index.
\end{definitionv}

\begin{definitionv}{synth_4}[Dyadic pilot resolution rule]
The pilot resolution \(K_0(n)\), for sample size \(n\in\mathbb N\), is defined by
\[
K_0(n)=
\begin{cases}
1, & n<n_0,\\[2pt]
2^{\left\lfloor \frac45\,\frac{\log m_0}{\log 2}\right\rfloor},
& n\ge n_0,
\end{cases}
\]
where \(n_0\) is the threshold for using the cubic estimator and \(m_0\) is the size of block \(0\) at sample size \(n\).
\end{definitionv}

\begin{definitionv}{synth_8}[Clipped marked-density pilot vector]
The pilot vector \(\widehat F(x)\), indexed by \(i\in\{0,\ldots,6\}\), is defined for real envelope parameters \(c_f,C_f\), a sample-size index \(n\in\mathbb N\), source and target samples \((\mathcal S_n,\mathcal T_n)\) each of size \(n\), and \(x\in\mathbb R\). For \(n<256\), set
\[
\widehat F(x)=\left(1,\frac12,\frac12,0,0,0,0\right).
\]
For \(n\ge256\), use \(K_0(n)\), the dyadic pilot resolution with exponent \(4/5\), and the block-\(0\) marked histogram \(H_{i,K_0(n)}^{(0)}(x)\) of \cref{obj:synth_7}, with the observable marks and zero-based record indexing of \cref{obj:synth_14}. Define \(\widehat F_i(x)\) by clipping this histogram value with the following lower and upper cutoffs, respectively:
\[
\begin{cases}
(c_f,C_f), & i=0,\\
(c_f/4,3C_f/4), & i\in\{1,2\},\\
(0,3C_f/4), & i\in\{3,4,5,6\}.
\end{cases}
\]
\end{definitionv}

\begin{definitionv}{synth_9}[Marked histogram residual]
The residual \(R_{i,K}^{(b)}(x)\) is defined by
\[
R_{i,K}^{(b)}(x)
=
H_{i,K}^{(b)}(x)-\widehat F_i(x),
\]
where the marked histogram \(H_{i,K}^{(b)}\) and the clipped pilot vector \(\widehat F\) are defined in \cref{obj:synth_7,obj:synth_8}, respectively. The parameters are arbitrary nonnegative integers \(n\) and \(K\), data consisting of \(n\) source records and \(n\) target covariates, a coordinate index \(i\) among the seven coordinates, a block index \(b\in\{0,1,2,3\}\), an evaluation point \(x\in\mathbb R\), and pilot-clipping parameters \(c_f,C_f\in\mathbb R\).
\end{definitionv}

\begin{definitionv}{synth_6}[Cubic correction resolution]
The cubic resolution \(K_3(n)\), for \(n\in\mathbb N\), is defined by
\[
K_3(n)=
\begin{cases}
1, & n<n_0,\\[2pt]
2^{\left\lfloor \dfrac{\log m_0}{\log 2}\right\rfloor}, & n\ge n_0,
\end{cases}
\]
where \(n_0\) is the threshold for using the cubic estimator and \(m_0\) is the size of block \(0\) at sample size \(n\). The floor is taken in the nonnegative integers, with value \(0\) when its argument is negative; the logarithm at \(0\) is assigned value \(0\).
\end{definitionv}

\begin{definitionv}{synth_5}[Quadratic correction resolution]
The quadratic resolution \(K_2(n)\), for \(n\in\mathbb N\), is defined by
\[
K_2(n)=
\begin{cases}
1, & n<n_0,\\[2pt]
2^{\left\lfloor \frac{4}{3}\,\frac{\log m_0}{\log 2}\right\rfloor},
& n\ge n_0,
\end{cases}
\]
where \(n_0\) is the threshold and \(m_0\) is the size of block \(0\) at sample size \(n\). The floor is taken in the nonnegative integers, with value \(0\) when its argument is negative.
\end{definitionv}

\begin{algorithmv}{def:cubic-estimator}[Cubic estimator \(\widehat T_A\)]
For \(A\in\{Y,D\}\), the inputs are the sample size \(n\), the threshold \(n_0\), the deterministic sample-channel blocks \(\mathcal B_b^G\) of sizes \(m_b\), the dyadic resolution maps \(K_0(n),K_2(n),K_3(n)\), the normalized marked histograms \(H_{i,K}^{(b)}\), the clipped pilot vector \(\widehat F\), the histogram residuals \(R_{i,K}^{(b)}\), the seven-coordinate index set \(\mathcal I\), the coordinate channels \(g(i)\), the observable marks \(W_{i,r}\), and the observed covariates \(X_r^{g(i)}\).
\begin{enumerate}
\item For \(n\ge n_0\), set the resolutions and cell midpoints to
\[
K_0=K_0(n),\qquad K_2=K_2(n),\qquad K_3=K_3(n),
\qquad
x_{\ell,K}=\frac{\ell-\tfrac12}{K},
\quad \ell=1,\ldots,K.
\]
The resolutions satisfy \(K_0\mid K_3\) and \(K_3\mid K_2\).

\item For \(n\ge n_0\), form the observable linear correction
\[
\mathsf L_A=
\sum_{i\in\mathcal I}
\left\{
\frac{1}{m_1}
\sum_{r\in\mathcal B_1^{g(i)}}
W_{i,r}\,\partial_i\Phi_A
\bigl(\widehat F(X_r^{g(i)})\bigr)
-
\int_0^1
\partial_i\Phi_A(\widehat F(x))\widehat F_i(x)\,dx
\right\}.
\]

\item For \(n\ge n_0\), form the observable quadratic correction
\[
\mathsf Q_A=
\frac{1}{2K_2}
\sum_{i,j\in\mathcal I}\sum_{\ell=1}^{K_2}
\partial_{ij}\Phi_A(\widehat F(x_{\ell,K_2}))
R_{i,K_2}^{(1)}(x_{\ell,K_2})
R_{j,K_2}^{(2)}(x_{\ell,K_2}).
\]

\item For \(n\ge n_0\), form the observable cubic correction
\[
\mathsf C_A=
\frac{1}{6K_3}
\sum_{i,j,k\in\mathcal I}\sum_{\ell=1}^{K_3}
\partial_{ijk}\Phi_A(\widehat F(x_{\ell,K_3}))
R_{i,K_3}^{(1)}(x_{\ell,K_3})
R_{j,K_3}^{(2)}(x_{\ell,K_3})
R_{k,K_3}^{(3)}(x_{\ell,K_3}).
\]
Repeated coordinate indices use distinct evaluation blocks.

\item Output
\[
\widehat T_A=
\begin{cases}
\displaystyle
\int_0^1\Phi_A(\widehat F(x))\,dx+
\mathsf L_A+\mathsf Q_A+\mathsf C_A,
& n\ge n_0,\\
0,& n<n_0.
\end{cases}
\]
Every integral is an exact finite sum, with \(K_0\mid K_3\) and \(K_3\mid K_2\). For \(n<n_0\), the score interval takes the finite-size fallback
\[
I_n=\Theta.
\]
\end{enumerate}
\end{algorithmv}

The estimator is an exact finite-sum cubic expansion around the clipped pilot. Distinct evaluation blocks supply repeated residual factors, while the nested dyadic partitions make the pilot-based derivative weights constant on correction cells.

\subsection{Observable reduction and mean-square control}

The following reduction identifies each marked coordinate with an observable measure and supplies the derivative envelope used in the Taylor remainder.

\begin{lemmav}{lem:observable-marked-reduction}[Observable marked density reduction]
Let \(n\in\mathbb N\) and let \(P\) be a transport law. Suppose:
\begin{itemize}
\item \textbf{(Lower density envelope.)} \(0<c_f<1\).
\item \textbf{(Upper density envelope.)} \(1<C_f\).
\item \textbf{(Hölder radius.)} \(1<L\).
\item \textbf{(Model membership.)} \(P\in\mathcal M_n\) for these constants, as specified in \cref{obj:def:model-class}.
\end{itemize}
Let \(F\), the seven-coordinate marked density vector, be as in \cref{obj:def:marked-density-vector}, with coordinates indexed by \(0,\ldots,6\). Write \(P_{\mathrm T}\) for the target covariate law and \(P_{\mathrm S}\) for the observed source-record law, whose record consists of covariate \(X\), encouragement \(Z\), receipt \(D\), and outcome \(Y\).

For every measurable \(B\subseteq[0,1]\) and every \(i\in\{0,\ldots,6\}\),
\[
\int_B F_i(x)\,dx
=
\begin{cases}
P_{\mathrm T}(B), & i=0,\\
\displaystyle\int_{\{X\in B\}}(1-Z)\,dP_{\mathrm S}, & i=1,\\
\displaystyle\int_{\{X\in B\}}Z\,dP_{\mathrm S}, & i=2,\\
\displaystyle\int_{\{X\in B\}}(1-Z)Y\,dP_{\mathrm S}, & i=3,\\
\displaystyle\int_{\{X\in B\}}ZY\,dP_{\mathrm S}, & i=4,\\
\displaystyle\int_{\{X\in B\}}(1-Z)D\,dP_{\mathrm S}, & i=5,\\
\displaystyle\int_{\{X\in B\}}ZD\,dP_{\mathrm S}, & i=6.
\end{cases}
\]
Moreover, the source arm densities \(q_0\) and \(q_1\) satisfy
\[
q_0(x)=F_1(x)\ge \frac{c_f}{4},
\qquad
q_1(x)=F_2(x)\ge \frac{c_f}{4},
\qquad x\in[0,1].
\]
For each \(A\in\{Y,D\}\), the transported contrast \(T_A\) and the integrand \(\Phi_A\) of \cref{obj:def:smooth-functional} satisfy
\[
T_A=\int_0^1\Phi_A(F(x))\,dx.
\]

For every \(A\in\{Y,D\}\) and every \(v\in\mathbb R^7\) in the deterministic clipping rectangle
\[
c_f\le v_0\le C_f,\qquad
\frac{c_f}{4}\le v_1,v_2\le\frac{3C_f}{4},\qquad
0\le v_i\le\frac{3C_f}{4}\quad(i=3,\ldots,6),
\]
every order \(k\in\{0,\ldots,4\}\) and every choice of coordinate indices \(i_1,\ldots,i_k\in\{0,\ldots,6\}\) obey
\[
\left|
\frac{\partial^k\Phi_A}{\partial v_{i_1}\cdots\partial v_{i_k}}(v)
\right|
\le
48(1+C_f)^2\left(\frac{4}{c_f}\right)^5.
\]
At order zero, the expression on the left is \(\lvert\Phi_A(v)\rvert\).
\end{lemmav}

The measure identities justify the normalized marked counts, and the lower arm-density bounds keep every functional denominator uniformly positive on the clipping rectangle. The next result combines pilot approximation, projection cancellation, and block-conditional variance calculations.

\begin{lemmav}{lem:marked-cubic-mse}[Uniform cubic mean-square bound]
Let \(c_f\) and \(C_f\) be the lower and upper density envelopes, and let \(L\) be the common Hölder radius. Suppose that:
\begin{itemize}
\item \textbf{(Envelope constants.)} \(0<c_f<1\), \(1<C_f\), and \(1<L\).
\item \textbf{(Sample size.)} The common source and target sample size \(n\) satisfies
\[
n\ge n_0=256.
\]
\item \textbf{(Model.)} \(P\in\mathcal M_n\), with \(\mathcal M_n\) defined in \cref{obj:def:model-class}.
\end{itemize}
For each \(A\in\{Y,D\}\), let \(\widehat T_A\) be the statistic in \cref{obj:def:cubic-estimator}, and let \(T_A\) be the target-averaged source contrast,
\[
T_A
=
\int_{\text{covariate space}}
f_{\mathrm T}(x)\bigl(m_{A1}(x)-m_{A0}(x)\bigr)\,dx.
\]
Then, under the equal-size source–target sampling law specified by \(P\),
\[
\mathbb E_P\!\left[(\widehat T_A-T_A)^2\right]
\le C_{\mathrm{mse}}\,n^{-2/3},
\]
where the finite, explicitly computable constant \(C_{\mathrm{mse}}=C_{\mathrm{mse}}(c_f,C_f,L)\) is
\[
\begin{aligned}
C_{\mathrm{mse}}
={}&3\Bigg\{
\left[2(1+C_f)^2(4/c_f)^5\right]^2 7^8
\Big[
2^7\,8!\,
\big\{(1+C_f+C_f^2)^4 5^{4/5}
 +(1+C_f+C_f^2)5^{7/5}\big\}\\
&\hspace{48mm}
+2^8[3(1+C_f)L]^8 5^{4/5}
\Big]\\
&\quad
+7^4[48(1+C_f)^2(4/c_f)^5]^2
 [3(1+C_f)L]^4 2^{1/2}5^{2/3}\\
&\quad
+2\Big[
\big\{24(1+C_f)^2(4/c_f)^5
       7^2[3(1+C_f)L]^2\big\}^2
 7^2\\
&\hspace{24mm}\cdot
\big\{2(1+C_f+C_f^2)5^{1/5}
 +2^{5/4}[3(1+C_f)L]^2 5^{1/5}\big\}
 2^{1/2}5^{1/2}\\
&\hspace{16mm}
+\big\{8(1+C_f)^2(4/c_f)^5
       7^3[3(1+C_f)L]^3\big\}^2
 2^{3/4}5^{3/4}
\Big]
\Bigg\}\\
&+3\Bigg\{
10\cdot7^2[48(1+C_f)^2(4/c_f)^5]^2\\
&\quad
+60\cdot7^4(2^2-1)^2
 [48(1+C_f)^2(4/c_f)^5]^2
 [1+2(1+C_f)]^4[2+C_f+C_f^2]^2\\
&\quad
+310\cdot7^6(2^3-1)^2
 [48(1+C_f)^2(4/c_f)^5]^2
 [1+2(1+C_f)]^6[2+C_f+C_f^2]^3
\Bigg\}.
\end{aligned}
\]
\end{lemmav}

The eighth-moment pilot bound controls the squared fourth-order remainder. The quadratic and cubic resolutions control projection errors, and conditioning on the training block preserves within-block coordinate dependence while retaining independence across the evaluation blocks. The displayed constant includes all approximation and fluctuation terms.

\subsection{Supporting length bounds}

The mean-square result calibrates both transported contrasts. Identification supplies the population affine-score identity at the target effect, and the bounded effect domain supplies the constant branch.

\begin{theoremv}{thm:honest-upper-full-elbow}[Honest upper length envelope]
Fix the noncoverage probability \(\alpha\), the density envelopes \(c_f,C_f\), and the Hölder radius \(L\), satisfying:
\begin{itemize}
\item \textbf{(Noncoverage.)} \(0<\alpha<1\).
\item \textbf{(Lower density envelope.)} \(0<c_f<1\).
\item \textbf{(Upper density envelope.)} \(1<C_f\).
\item \textbf{(Hölder radius.)} \(1<L\).
\end{itemize}
The score interval \(I_n\) of \cref{obj:def:score-interval} belongs to the class \(\mathfrak H_n\) of honest intervals in \cref{obj:def:honest-intervals} for every positive integer \(n\).

Set the sample-size threshold to
\[
n_0=256.
\]
Write \(\mu(P)\) for the transported first stage, namely the target-averaged source receipt contrast:
\[
\mu(P)=T_D(P).
\]
There exists a finite constant \(C_0=C_0(\alpha,c_f,C_f,L)>0\) such that, for every integer \(n\ge n_0\) and every \(P\in\mathcal M_n\), the model class of \cref{obj:def:model-class},
\[
\mathbb E_P\!\left[\lambda_{\Theta}(I_n)\right]
\le C_0\min\left\{1,\frac{n^{-1/3}}{\mu(P)}\right\},
\]
where \(\lambda_{\Theta}\) denotes Lebesgue length restricted to the effect domain \(\Theta\), and the expectation is under the joint law of the source and target samples, each of size \(n\).

With the same constant \(C_0\), for every integer \(n\ge n_0\) and every real \(a\) satisfying \(0<a\le 1/4\), the minimax expected-length frontier \(L_n(a)\) of \cref{obj:def:length-frontier} satisfies
\[
L_n(a)\le C_0\min\left\{1,\frac{n^{-1/3}}{a}\right\}.
\]
\end{theoremv}

The upper result gives all-positive-sample-size honesty and, from \(n_0\) onward, both a law-specific length bound and a uniform bound on each strength slice. The converse is stated next; its causally legal experiment and testing argument are developed in \cref{sec:lower-appendix}.

\begin{theoremv}{thm:honest-lower-full-elbow}[Honest length lower bound]
Fix constants satisfying:
\begin{itemize}
\item \textbf{(Coverage.)} The noncoverage probability satisfies \(0<\alpha<1\).
\item \textbf{(Density envelopes.)} The lower and upper density envelopes satisfy \(0<c_f<1<C_f\).
\item \textbf{(Smoothness.)} The common Hölder radius satisfies \(L>1\).
\end{itemize}
Let \(\mathcal M_n\) be the model class in \cref{obj:def:model-class}, and let \(\mathcal M_n(a)\) be its strength slice in \cref{obj:def:strength-slice}. Write \(\theta(P)\) for the target-population average treatment effect among compliers under \(P\), and let the effect domain and sample-size threshold be
\[
\Theta=[-1,1],
\qquad n_0=256.
\]
There exists a constant \(c_0=c_0(\alpha,c_f,C_f,L)>0\) such that, for every integer \(n\ge n_0\) and every \(0<a\le 1/4\), the following holds. Every measurable confidence-set procedure \(C_n\), based on \(n\) source records and \(n\) target covariate records, with \(C_n(\omega)\subseteq\Theta\) for every sample realization \(\omega\), and satisfying
\[
P\bigl(\theta(P)\in C_n\bigr)\ge 1-\alpha
\qquad\text{for every }P\in\mathcal M_n,
\]
obeys
\[
\sup_{P\in\mathcal M_n(a)}
\mathbb E_P\!\left[\lambda_{\Theta}(C_n)\right]
\ge c_0\min\left\{1,\frac{n^{-1/3}}{a}\right\},
\qquad
\lambda_{\Theta}(C_n)
=\operatorname{Leb}(C_n\cap\Theta).
\]
Here probabilities and expectations are taken under the observed two-sample law. Moreover, the minimax expected-length frontier \(L_n(a)\) of \cref{obj:def:length-frontier} satisfies
\[
L_n(a)\ge c_0\min\left\{1,\frac{n^{-1/3}}{a}\right\}.
\]
\end{theoremv}

Together these bounds yield the exact frontier stated in \cref{obj:thm:sharp-length-frontier-full-elbow}. Full proof scripts for the anchored statements are collected in \cref{sec:deferred-proofs}.
